%%=================================================================
%% CHAPTER 1: Secluded Majorana SIDM — Phenomenology
%% Source: github.com/0mp3s/Secluded-Majorana-SIDM (DOI: 10.5281/zenodo.19225823)
%%=================================================================

\section{Introduction}
\label{ch1:intro}

The small-scale structure anomalies of $\Lambda$CDM — the core-cusp problem,
missing satellites, too-big-to-fail — motivate self-interacting dark matter (SIDM),
where DM particles scatter elastically via a light mediator. We consider a
\emph{secluded} dark sector: a Majorana fermion $\chi$ coupled to a real scalar
$\phi$ with no tree-level coupling to the SM.

\section{The Model}
\label{ch1:model}

The interaction Lagrangian is:
\begin{equation}
  \Lc_{\rm int} = -\frac{y}{2}\bar{\chi}\chi\phi
\end{equation}

In the non-relativistic limit, scalar exchange generates an attractive Yukawa potential:
\begin{equation}
  V(r) = -\frac{\alpha}{r}e^{-\mphi r}, \qquad \alpha = \frac{y^2}{4\pi}
\end{equation}

The model has three free parameters: $(m_\chi,\, m_\phi,\, \alpha)$.

\subsection{Scalar Potential and Vacuum Stability}

The full scalar potential:
\begin{equation}
  V(\phi) = \frac{1}{2}m_\phi^2\phi^2 + \frac{\mu_3}{3!}\phi^3 + \frac{\lambda_\phi}{4!}\phi^4
\end{equation}

Vacuum stability (no secondary minimum deeper than $\phi=0$) requires:
\begin{equation}
  \lambda_\phi > \frac{3\mu_3^2}{8m_\phi^2}
\end{equation}

For the cannibal lower bound $\mu_3/m_\phi = 1.7$: $\lambda_\phi \gtrsim 1.08$.
Numerically: $\lambda_\phi \gtrsim 0.96$ (from $V(\phi_{\rm min}) > 0$ scan).

\section{Self-Interaction Cross Section: Variable Phase Method}
\label{ch1:vpm}

The transfer cross section is computed via the VPM ODE for phase shifts $\delta_\ell$:
\begin{equation}
  \frac{d\delta_\ell}{dx} = +\frac{\lambda\,e^{-x}}{\kappa\,x}
  \left[\hat{j}_\ell(\kappa x)\cos\delta_\ell - \hat{n}_\ell(\kappa x)\sin\delta_\ell\right]^2
\end{equation}

where $x = m_\phi r$, $\kappa = \mu v_{\rm rel}/m_\phi$, $\lambda = \alpha m_\chi/m_\phi$.

For identical Majorana fermions, quantum statistics requires:
\begin{equation}
  \sigmaT = \frac{2\pi}{k^2}\left[\sum_{l\,\rm even}(2l+1)\sin^2\delta_l
             + 3\sum_{l\,\rm odd}(2l+1)\sin^2\delta_l\right]
\end{equation}

\textbf{Key result (proved in Chapter~2, Appendix~\ref{app:vpm_proof}):} VPM is
the \emph{exact} fluctuation determinant of the path integral — no approximation.

\section{Relic Density}
\label{ch1:relic}

Freeze-out via $\chi\chi \to \phi\phi$ (t/u-channel scalar mediator).
For Majorana fermions, this is s-wave:
\begin{equation}
  \langle\sigma v\rangle_0 = \frac{\pi\alpha^2}{4m_\chi^2}
\end{equation}

Sommerfeld enhancement at freeze-out: $S_0 < 1.026$ for all 17 relic-viable BPs
(correction $< 2.5\%$). CMB constraint: inapplicable (secluded sector, $f_{\rm eff}=0$).

\section{MCMC Posterior and Viable Island}
\label{ch1:mcmc}

A Bayesian MCMC scan (32 walkers, 5000 steps, emcee, fast JIT solver) over
$(m_\chi, m_\phi, \alpha)$ against 9 astrophysical observables identifies
the viable island:
\begin{align}
  m_\chi &\in [10,\,100]\,{\rm GeV} \\
  m_\phi &\in [7.6,\,14.8]\,{\rm MeV} \\
  \alpha  &\in [5\times10^{-4},\,5\times10^{-3}]
\end{align}

122 benchmark points pass all cuts simultaneously.

\section{Astrophysical Confrontation}
\label{ch1:astro}

[Section to be expanded. Key results summarized in Table~\ref{tab:ch1_pheno}.]

\begin{table}[h!]
\centering
\caption{Phenomenological summary for named benchmark points.}
\label{tab:ch1_pheno}
\begin{tabular}{lcccc}
\toprule
Observable & BP1 & BP9 & MAP & Constraint \\
\midrule
Fornax GC survival & 12/15 & 13/15 & \textbf{14/15} & all 5 survive \\
$r_{\rm core}$ Fornax [pc] & 449 & 541 & \textbf{829} & $\gtrsim 500$ \\
Fornax $\chi^2$ (Jeans+OM) & 35 & — & \textbf{2.3} & — \\
Read+2019 bound & \checkmark & \textsf{FAIL} & $\checkmark$ (margin +0.05) & $\sigma/m < 1.5$ cm$^2$/g \\
RAR scatter & \textbf{0.122} dex & — & 0.28 dex & 0.177 dex (obs.) \\
Sommerfeld $S_0$ at $T_{fo}$ & 1.003 & 1.010 & 1.025 & $< 2.5\%$ \\
\bottomrule
\end{tabular}
\end{table}
